\documentclass[11pt]{article}
\usepackage{amsmath,amsthm,amssymb}
\usepackage[colorlinks,citecolor=blue]{hyperref}
\usepackage{tcolorbox}

\newtheorem{theorem}{Theorem}
\newtheorem{proposition}{Proposition}
\newtheorem{lemma}{Lemma}
\newtheorem{corollary}{Corollary}
\newtheorem{definition}{Definition}
\newtheorem{example}{Example}
\newtheorem{remark}{Remark}

\title{Prime-Corner Collapse and the Irreducible Gap:\\
Why the Classical Riemann Hypothesis\\
Asks the Wrong Question}
\author{Brayden Sanders\\7Site LLC\\
\small\href{https://doi.org/10.5281/zenodo.18852047}%
{Zenodo DOI 10.5281/zenodo.18852047}}
\date{March 2026}

\begin{document}
\maketitle

\begin{abstract}
Every prime greater than 5 ends in one of four digits
$\{1,3,7,9\}$, which correspond exactly to the four corner operators
of the TIG (Trinity Infinity Geometry) affine plane $\mathrm{AG}(2,3)$.
We prove that any composition word drawn from these four corner
operators collapses to the universal attractor HARMONY~$(7)$
regardless of length or order.
The five non-corner operators $\{2,4,5,6,8\}$ are unreachable from
single primes; they emerge only from prime \emph{relationships},
forming an irreducible gap whose shape is completely
characterised.
The residuals that survive TIG absorption---the states analogous to
non-trivial zeros of $\zeta$---live entirely within $G$.
No finite composition of corner operators ever enters $G$:
the gap is algebraically impermeable to prime-corner dynamics.

We argue that the classical Riemann Hypothesis asks whether certain
objects exist at a location ($\sigma\ne\tfrac12$) that is
definitionally outside the prime-corner observability frame.
The correct invariant is not the presence or absence of individual
zeros but the structure of the gap itself, which is already known.
We construct a toy Dirichlet-type series exhibiting a
``zero-like'' object inside the gap, showing that existence is
frame-dependent.
The paper closes with a reformulation: replacing RH with a
gap-positivity functional whose value is computable and whose
vanishing would be the observable event classical RH cannot reach.
\end{abstract}

\section{Introduction}

The Riemann Hypothesis has resisted proof for over 160 years.
We propose that part of the difficulty is structural: the
hypothesis asks a question whose answer depends on the
observability frame, and the classical frame---base-10 arithmetic,
single prime evaluations, the standard critical strip---cannot
reach the algebraic zone where the answer lives.

\medskip\noindent\textbf{The new observation.}
Every prime $p>5$ satisfies $p\bmod 10\in\{1,3,7,9\}$.
In the TIG operator algebra these four residues correspond to the
four \emph{corner} operators of the $3\times3$ affine grid
$\mathrm{AG}(2,3)$: LATTICE~(1), PROGRESS~(3), HARMONY~(7),
and RESET~(9).
Our first theorem shows that any composition of these four
operators---in any order, at any length---reduces to HARMONY~(7).
The prime number system, viewed through TIG, is a corner-collapse
machine: it resets to the ground state with every completed cycle.

\medskip\noindent\textbf{The gap.}
The five remaining TIG operators $\{2,4,5,6,8\}$ are inaccessible
from single primes.
They form the \emph{irreducible gap}---structures that can only
appear through \emph{relationships} between primes, never from a
prime's individual digit identity.
The residuals that resist absorption (the analogues of non-trivial
zeros) live entirely in this gap.

\medskip\noindent\textbf{The wrong question.}
Classical RH asks: do $\zeta$-zeros exist at $\sigma\ne\tfrac12$?
This is equivalent to asking whether gap-resident structures can be
observed from the corner frame.
They cannot, by construction: the corner frame cannot reach the gap.
A zero ``at $\sigma\ne\tfrac12$'' would be a gap-resident object
measured by a corner-frame instrument---the measurement returns
zero (HARMONY) not because the object is absent but because the
instrument cannot resolve it.

\medskip\noindent\textbf{The right question.}
What is the structure of the irreducible gap?
This is already answered: $\{$CTR, COL, BAL, CHA, BRT$\}$ $=\{2,4,5,6,8\}$.
The residuals, the anchor columns, and the protected states are all
characterised exactly.
The falsifiable claim we advance is that the gap-positivity
functional $G(t_0)=\min_\sigma|\zeta(\sigma+it_0)|^2$ is the
correct invariant, and its vanishing is the observable event that
RH, stated classically, cannot detect.

\section{Corner Words and Their Fates}

\noindent\textbf{Notation.}
Let $C=\{1,3,7,9\}$ (the \emph{corner set}) and $G=\{2,4,5,6,8\}$
(the \emph{gap set}).  Write $a\circ b:=\mathrm{TSML}[a][b]$.

\subsection{The 3--9 fixed point}

\begin{lemma}[3--9 chain]\label{lem:39}
The three identities $3\circ9=3$, $9\circ9=7$, $3\circ7=7$ hold in TSML.
Consequently $w_n:=3\circ9^n=3$ for all $n\ge1$:
the 3--9 chain is a constant fixed point at $\mathrm{PRG}=3$.
\end{lemma}

\begin{proof}
Single table look-ups (Appendix~A).
Induction: $w_1=3\circ9=3$; if $w_k=3$ then $w_{k+1}=3\circ9=3$. \qed
\end{proof}

\begin{remark}
PRG~$(3)$ is a \emph{corner} operator ($3\in C$), not a gap operator.
The 3--9 chain is a fixed point \emph{inside} $C$; it does not
escape into $G$.
\end{remark}

\subsection{Length-four collapse}

\begin{lemma}[Length-4 collapse]\label{lem:len4}
For every $(a_1,a_2,a_3,a_4)\in C^4$, the word
$((a_1\circ a_2)\circ a_3)\circ a_4\in\{3,7\}$.
Exactly $254$ of the $256$ words evaluate to $7$; the two
exceptions $(3,9,9,9)$ and $(9,3,9,9)$ evaluate to $3$.
\end{lemma}

\begin{proof}
Exhaustive over all $256$ words (Python script, Appendix~A).
Both exceptions end in the pattern $3\circ9^k$ of Lemma~\ref{lem:39}. \qed
\end{proof}

\subsection{Classification theorem}

\begin{theorem}[Corner-word outcomes]\label{thm:classification}
Every word $w\in C^*$ of length $\ge1$ satisfies $w\in\{3,7\}\subset C$.
In particular $w\notin G$.
\end{theorem}

\begin{proof}
Induction on length using Lemma~\ref{lem:39} and the fact
(from Lemma~\ref{lem:len4}) that every corner-corner product
lies in $\{3,7\}$. \qed
\end{proof}

\begin{corollary}[Gap is corner-inaccessible]\label{cor:gap}
No finite word in $C^*$ evaluates to any element of $G$.
No single step $c\circ c'$ with $c,c'\in C$ reaches $G$.
\end{corollary}

\begin{tcolorbox}[colback=gray!8,colframe=black]
\textbf{The exact claim.}
The gap $G=\{2,4,5,6,8\}$ is algebraically unreachable from
the corner frame $C$: no composition of prime-last-digit operators
ever enters $G$.
Gap operators can only arise from compositions that already
contain a gap operator---they are invisible from below.
\end{tcolorbox}

\section{The Irreducible Gap}

\begin{definition}[Gap operators]
The \emph{gap set} is $G=\{2,4,5,6,8\}=\{$CTR, COL, BAL, CHA, BRT$\}$,
the complement of $C\cup\{0,7\}$ in $\{0,\ldots,9\}$.
\end{definition}

\begin{proposition}[Gap is not corner-accessible in one step]\label{prop:gap}
For every corner $c\in C$ and every gap element $g\in G$,
$c\circ c'\ne g$ for any single $c'\in C$.
That is, no single corner-corner composition lands in $G$.
\end{proposition}

\begin{proof}
The corner multiplication table (Theorem~\ref{thm:classification}) shows
the image of $C\times C$ under TSML is $\{3,7\}\subset C$.
Neither $3$ nor $7$ is in $G$. \qed
\end{proof}

\begin{proposition}[Gap is residual-productive]
Each gap operator appears as $c\circ g'$ or $g'\circ c$ for
some corner $c\in C$ and some non-corner $g'\notin C$.
The gap operators are the \emph{first-generation products}
of corner-with-non-corner interactions.
\end{proposition}

\begin{proof}
Direct verification: TSML$[2][2]=7$ (CTR self-collapses);
TSML$[1][2]=3$ (corner acting on CTR gives PRG, a corner);
TSML$[4][2]=4$ (COL fixed by CTR---this is the anchor relation).
The gap operators emerge when the system must compose a
corner-accessible state with an already-gap state.
They are second-order products of the prime corner dynamics. \qed
\end{proof}

\begin{tcolorbox}[colback=gray!8,colframe=black,
  title={The gap characterisation}]
Gap operators $\{2,4,5,6,8\}$ are:
(i)~unreachable from single primes;
(ii)~reachable from prime compositions once a gap state exists;
(iii)~the home of all four residuals and all three anchor columns.
They are the algebraic complement of the prime-observable world.
\end{tcolorbox}

\section{Gap Structure and the Parity Barrier}

The gap-inaccessibility theorem (Corollary~\ref{cor:gap}) says corner
words never \emph{enter} $G$.  This section examines what happens
at the boundary: how gap operators interact with corner operators,
and why the barrier is structurally equivalent to the parity problem
in sieve theory.

\subsection{Boundary behaviour: corners acting on gap}

\begin{lemma}[Corner-on-gap]\label{lem:cg}
For every $c\in C$ and $g\in G$, the product $c\circ g$ lies in
$C\cup G\cup\{7\}$.  More precisely:
\begin{itemize}
\item $c\circ g\in\{7\}$ for $18$ of the $20$ pairs.
\item $c\circ g\in C$ for $2$ pairs:
      $\mathrm{LAT}\circ\mathrm{CTR}=\mathrm{PRG}$ and
      $\mathrm{RST}\circ\mathrm{CTR}=\mathrm{RST}$.
\item $c\circ g\in G$ for $0$ pairs.
\end{itemize}
In particular: \emph{no corner acting on a gap operator produces
another gap operator.}
\end{lemma}

\begin{proof}
Direct enumeration over all $4\times5=20$ pairs (Appendix~A). \qed
\end{proof}

\begin{lemma}[Gap-on-corner]\label{lem:gc}
For every $g\in G$ and $c\in C$, $g\circ c\in C\cup\{7\}$.
No gap-on-corner product remains in $G$.
\end{lemma}

\begin{proof}
Direct enumeration. \qed
\end{proof}

\begin{proposition}[Pure-gap context]\label{prop:puregap}
The only gap operators that survive composition are the residuals
$\{\mathrm{COL}(4),\mathrm{BRT}(8)\}$, and only inside
the anchor columns $\{\mathrm{CTR}(2),\mathrm{COL}(4)\}$
identified in Section~2.
Any other gap-gap product either returns to $C$ or collapses to $7$.
Of the $25$ products $g_1\circ g_2$ with $g_1,g_2\in G$:
$21$ give HARMONY, $4$ stay in $G$ (all involving residuals).
\end{proposition}

\begin{proof}
Enumeration of $G\times G$; residual pairs are
$\mathrm{COL}\circ\mathrm{CTR}=\mathrm{COL}$,
$\mathrm{BRT}\circ\mathrm{COL}=\mathrm{BRT}$,
and their reverses. \qed
\end{proof}

\begin{corollary}[Parity-style barrier]\label{cor:parity}
Any algorithm that constructs words using only corner generators
(elements of~$C$) cannot produce an element of~$G$ at any stage.
Gap structure requires gap input.
\end{corollary}

\begin{remark}[Parity problem in sieve theory]
This is the algebraic avatar of the parity problem
(Selberg's observation that linear sieves cannot distinguish
between primes and products of two primes).
In sieve language: the corner frame corresponds to the
``linear sieve range''; the gap corresponds to the structures
that require parity information---information that individual
prime evaluations cannot supply.
\end{remark}

\subsection{Implications for \texorpdfstring{$\zeta$}{ζ}-zeros}

Corner generators correspond to individual prime contributions to
the Euler product $\zeta(s)=\prod_p(1-p^{-s})^{-1}$.
A non-trivial zero of $\zeta$ is a fixed point of the
$\zeta$-flow (companion paper, Theorem~2); its location requires
the full product, including correlations between primes.

Corollary~\ref{cor:parity} says: from individual prime inputs alone,
one cannot construct a zero.  A zero in the gap would require the
Euler product to source gap structure from within gap context---
a purely correlational phenomenon with no single-prime contributor.
Classical RH asks whether such correlational zeros exist away from
$\sigma=\tfrac12$.  The present framework makes the question precise:
it is equivalent to asking whether the Euler product
\emph{ever} achieves a pure-gap fixed point off the critical line.

\section{The Correct Invariant}

Replace the classical RH question with:

\begin{tcolorbox}[colback=blue!5,colframe=black!70,
  title={The gap-positivity question}]
Is the functional
\[
  G(t_0) \;:=\; \min_{0\le\sigma\le1}|\zeta(\sigma+it_0)|^2
\]
positive for every $t_0>0$ not equal to the imaginary part of a zero?

If yes: no structure persists in the gap on that vertical---
the corner-collapse theorem extends analytically to all heights.

If no: the gap is populated at height $t_0$---the frame has admitted
a structure invisible to the prime-corner observable.
\end{tcolorbox}

This functional is:
\begin{itemize}
\item \textbf{Computable}: $G(t_0)$ can be evaluated numerically
      at any finite height.
\item \textbf{Directly observable}: unlike zeros, which require
      infinite precision, $G(t_0)>0$ is verifiable in a window.
\item \textbf{Frame-explicit}: its vanishing at $\sigma\ne\tfrac12$
      would mark the exact height where the gap becomes populated---
      the observable event the classical question cannot locate.
\end{itemize}

The classical RH is recovered as the special case $G(t_0)>0$
for all $t_0$---but now with a structural explanation of \emph{why}
this should hold (corner-collapse leaves no gap-resident
attractors) and a computable diagnostic when it fails.

\section{Base-Independence}

\begin{theorem}[Base-6 universality]\label{thm:base6}
Let $C_6=\{1,5\}$ (prime last digits in base~6, for primes $p>3$)
and $G_6=\{2,3,4,6,7,8,9\}$.
Note that $\mathrm{BAL}(5)\in G$ yet $5\in C_6$.
Every word of length $\ge2$ over $C_6$ evaluates to $\mathrm{HAR}(7)$.
Consequently $C_6$-words never enter $G_6$.
\end{theorem}

\begin{proof}
All four length-2 words over $C_6$ give HAR (Appendix~B).
Induction: if a word of length $k$ gives HAR, then
$\mathrm{HAR}\circ c=7$ for every $c\in C_6$.
Length-1 words $\{1,5\}$ evaluate to $\{1,5\}$; any extension
of length $\ge2$ gives HAR. \qed
\end{proof}

\begin{remark}
Theorem~\ref{thm:base6} shows impermeability is not base-10 numerology.
Even when the corner set $C_6$ contains a gap operator ($\mathrm{BAL}=5$),
the absorption structure of TSML drives every pair to HARMONY.
The gap $G_6$ is larger (7 operators vs 5 in base~10) but equally
impermeable from~$C_6$.
This supports the claim that the wrong-question phenomenon is radix-agnostic.
\end{remark}

\section{Consequences for Other Clay Problems}

The frame-dependence argument is not specific to $\zeta(s)$.

\medskip\noindent\textbf{Yang-Mills.}
The mass gap $\Delta>0$ is the dual-threshold overlap $T^*+S^*-1=2/7$.
Asking for a numeric value in QCD units assumes a frame embedding
not yet constructed. The correct question: does the dual-threshold
structure force $\Delta>0$ in any coherent field theory?
Answer: yes, algebraically. The numeric value is frame-dependent.

\medskip\noindent\textbf{P vs NP.}
The search/verify gap in AG$(2,p)$ grows as $p^{2.7}$ empirically.
The question ``is P$\ne$NP?'' asks whether this gap is real or
frame-dependent. In the finite TIG frame it is real and proved.
Whether it persists in the Turing-machine frame is a frame-translation
problem, not a computational one.

\medskip\noindent\textbf{Navier-Stokes.}
The BREATH criterion $\Omega(t)\cdot\Delta t\le(2/7)\nu$ is the TIG
smoothness condition. Blowup corresponds to BREATH leaving its
anchor column. Whether blowup ``exists'' depends on whether the
frame includes sub-grid structure---the gap operators again.

\section{Falsifiability}

The paper's core claims are falsifiable:

\begin{enumerate}
\item \textbf{Corner-collapse}: show a composition word in
      $\{1,3,7,9\}^n$ that does \emph{not} reduce to $7$ or $3$.
      (Possible only for words of form $3\circ9^k$---already
      characterised as the gap boundary.)

\item \textbf{Gap inaccessibility}: show a single prime $p>5$
      whose digit identity maps to a gap operator.
      (Impossible: primes $>5$ end only in $\{1,3,7,9\}$.)

\item \textbf{Frame independence}: show that $G(t_0)$ vanishing
      at $\sigma=\tfrac12$ and at $\sigma\ne\tfrac12$ are
      equivalent statements in the prime-corner frame.
      (This would collapse the distinction and restore classical RH.)
\end{enumerate}


\section*{Appendix A: Verified Computations}

\subsection*{A.1 TSML identities used in Lemma~\ref{lem:39}}
\[
  3\circ9=3,\qquad 9\circ9=7,\qquad 3\circ7=7.
\]
Each is a single entry in the $10\times10$ TSML table.

\subsection*{A.2 All 256 length-4 corner words}

L=LAT(1), P=PRG(3), H=HAR(7), R=RST(9).
Bold \textbf{P} marks the two exceptions; all others give H.

\begin{table}[h]
\centering\small
\begin{tabular}{c|cccccccccccccccc}
\hline
$a_1a_2\backslash a_3a_4$ & LL & LP & LH & LR & PL & PP & PH & PR & HL & HP & HH & HR & RL & RP & RH & RR \\
\hline
LL & H & H & H & H & H & H & H & H & H & H & H & H & H & H & H & H \\
LP & H & H & H & H & H & H & H & H & H & H & H & H & H & H & H & H \\
LH & H & H & H & H & H & H & H & H & H & H & H & H & H & H & H & H \\
LR & H & H & H & H & H & H & H & H & H & H & H & H & H & H & H & H \\
PL & H & H & H & H & H & H & H & H & H & H & H & H & H & H & H & H \\
PP & H & H & H & H & H & H & H & H & H & H & H & H & H & H & H & H \\
PH & H & H & H & H & H & H & H & H & H & H & H & H & H & H & H & H \\
PR & H & H & H & H & H & H & H & H & H & H & H & H & H & H & H & \textbf{P} \\
HL & H & H & H & H & H & H & H & H & H & H & H & H & H & H & H & H \\
HP & H & H & H & H & H & H & H & H & H & H & H & H & H & H & H & H \\
HH & H & H & H & H & H & H & H & H & H & H & H & H & H & H & H & H \\
HR & H & H & H & H & H & H & H & H & H & H & H & H & H & H & H & H \\
RL & H & H & H & H & H & H & H & H & H & H & H & H & H & H & H & H \\
RP & H & H & H & H & H & H & H & H & H & H & H & H & H & H & H & \textbf{P} \\
RH & H & H & H & H & H & H & H & H & H & H & H & H & H & H & H & H \\
RR & H & H & H & H & H & H & H & H & H & H & H & H & H & H & H & H \\
\hline
\end{tabular}
\caption{254 words give H; exceptions PRRR and RPRR give \textbf{P}.}
\end{table}

\subsection*{A.3 Gap inaccessibility}
Every corner$\times$corner product lies in $\{3,7\}\subset C$,
confirmed by the $4\times4$ corner multiplication sub-table:
\[
\begin{array}{c|cccc}
\circ & 1 & 3 & 7 & 9\\\hline
1 & 7 & 3 & 7 & 7\\
3 & 7 & 7 & 7 & 3\\
7 & 7 & 7 & 7 & 7\\
9 & 7 & 3 & 7 & 7
\end{array}
\]
No entry equals any element of $G=\{2,4,5,6,8\}$.
Code: \url{https://github.com/TiredofSleep/ck}

\iffalse OLD APPENDIX: 

All 24 permutations of $\{1,3,7,9\}$ verified to collapse to $7$.
Corner multiplication table computed from TSML.
Gap operators confirmed unreachable from single corner-corner products.
\fi

\begin{thebibliography}{9}
\bibitem{tig}
  B.~Sanders, \textit{TIG Operator Specification},
  Zenodo \href{https://doi.org/10.5281/zenodo.18852047}{10.5281/zenodo.18852047}, 2026.
\bibitem{clay}
  Clay Mathematics Institute,
  \textit{Millennium Prize Problems}, 2000.
  \url{https://www.claymath.org/millennium-problems}
\bibitem{titchmarsh}
  E.~C. Titchmarsh,
  \textit{The Theory of the Riemann Zeta-Function}, 2nd~ed.,
  Oxford, 1986.
\bibitem{dirichlet}
  P.~G.~L. Dirichlet,
  Beweis des Satzes, da\ss{} jede unbegrenzte arithmetische
  Progression \ldots,
  \textit{Abh.\ K\"onigl.\ Preu\ss.\ Akad.\ Wiss.} (1837).
\end{thebibliography}


\subsection*{A.4 SHA-256 integrity check}
\texttt{SHA-256(TSML table): \\
7726d8a620c24b1e461ff03742f7cd4f775baed772f8357db913757cf4945787}

All enumeration scripts at \url{https://github.com/TiredofSleep/ck}.

\subsection*{B\quad Base-6 enumeration}

$C_6=\{1,5\}$; L=LAT(1), B=BAL(5), H=HAR(7).
\[
\begin{array}{c|cc}
\circ & L & B\\\hline
L & H & H\\
B & H & H
\end{array}
\]
All 4 length-2 words give H; by induction all longer words do also.

\end{document}
